% !TeX root = main.tex
\section{Conclusion and Future Work}

This work presented the graphics-based Parallel Collision Detection
(gPCD) framework, a GPU-driven collision-detection architecture that
integrates broad-phase occupancy generation directly into the graphics
pipeline. By allowing occupancy construction and rendering to proceed
concurrently, the framework introduces a new execution architecture for
GPU-based collision detection while preserving the conventional
broad-phase / narrow-phase organization.

Experimental results demonstrated near-linear scaling across the
evaluated particle range and supported the complexity analysis
developed in this work. The framework successfully processed more than
11 million particles on a single CPU-GPU system while maintaining high
throughput and interactive performance. Locality studies further showed
that memory organization influences implementation efficiency without
altering the underlying asymptotic behavior of the framework.

These results suggest that graphics-pipeline occupancy generation is a
viable approach for large-scale particle collision detection and that a
single CPU-GPU system can support particle-system scales traditionally
associated with high-performance computing environments.

Future work will investigate larger GPU configurations, multi-GPU and
distributed execution, and the application of graphics-pipeline
occupancy generation to particle-based thermo-fluid and multiphysics
simulations.

